\appendixheading

*Lemma 1. --- Let X be a locally compact space, R an open equivalence relation in X, such that the quotient space X/R is paracompact; let \(\pi\) be the canonical mapping of X onto X/R. There exists a continuous real-valued function \(F \geqslant 0\) on X such that:

\begin{enumerate}
\def\labelenumi{\alph{enumi})}
\item
  F is not identically zero on any equivalence class with respect to R;
\item
  for every compact subset K of X/R, the intersection of \(\overline{\pi}^{1}(K)\) with Supp F is compact.
\end{enumerate}

To each point \(u \in \mathbf{X} / \mathbf{R}\) let us associate a function \(f_{u} \in \mathcal{K}_{+}(\mathbf{X})\) such that \(f_{u}\) is not identically zero on \(\overline{\pi}(u)\) ; let \(\Omega_{u}\) be the open set of points where \(f_{u} > 0\) ; thus \(u \in \pi(\Omega_{u})\) . Since \(\pi\) is an open mapping, the \(\pi(\Omega_{u})\) form an open covering of \(\mathbf{X} / \mathbf{R}\) . There exists a locally finite open covering \((\mathrm{U}_{\iota})_{\iota \in \mathbf{I}}\) , finer than the covering by the \(\pi(\Omega_{u})\) , then (GT, IX, §4, No.~3, Prop. 3) a partition of unity \((g_{\iota})_{\iota \in \mathbf{I}}\) on \(\mathbf{X} / \mathbf{R}\) subordinate to the covering \((\mathrm{U}_{\iota})\) . For every \(\iota \in \mathbf{I}\) , choose a \(u_{\iota}\) such that \(\mathrm{U}_{\iota} \subset \pi(\Omega_{u_{\iota}})\) . The function \(\mathrm{F}_{\iota} = (g_{\iota} \circ \pi) \cdot f_{u_{\iota}}\) belongs to \(\mathcal{K}(\mathbf{X})\) and has support contained in \(\overline{\pi}(\mathrm{U}_{\iota})\) . The supports of the \(F_{\iota}\) therefore form a locally finite family, so that one can define a continuous function \(F \geqslant 0\) on \(X\) by setting \(F = \sum_{\iota \in I} F_{\iota}\) .

For every \(u \in X/R\) , there exists an \(\iota\) such that \(g_{\iota}(u) > 0\) and therefore \(u \in U_{\iota}\) ; next, there exists an \(x \in \Omega_{u_{\iota}}\) such that \(\pi(x) = u\) ; then \(f_{u_{\iota}}(x) > 0\) and \(g_{\iota}(\pi(x)) > 0\) , therefore \(F_{\iota}(x) > 0\) and a fortiori \(F(x) > 0\) ; this proves that F has the property a). Finally, let K be a compact subset of X/R. There exists a finite subset J of I such that, for \(\iota \in I - J\) , one has \(\mathrm{U}_{\iota}\cap \mathrm{K} = \varnothing\) , therefore \(\overline{\pi} (\mathbf{K})\cap \operatorname {Supp}\mathbf{F}_{\iota} = \varnothing\) . Then the set

\[
\overline {{\pi}} ^ {1} (\mathrm{K}) \cap \operatorname{Supp} \mathrm{F} = \overline {{\pi}} ^ {1} (\mathrm{K}) \cap \left(\bigcup_ {\iota \in \mathrm{I}} \operatorname{Supp} \mathrm{F} _ {\iota}\right) = \overline {{\pi}} ^ {1} (\mathrm{K}) \cap \left(\bigcup_ {\iota \in \mathrm{J}} \operatorname{Supp} \mathrm{F} _ {\iota}\right)
\]

is compact.

Lemma 2. --- Let G be a locally compact group countable at infinity, M a Baire space. Suppose that G operates on the left continuously and transitively in M. For every \(x \in M\) , let \(H_{x}\) be the stabilizer of x in G, so that the mapping \(s \mapsto sx\) of G onto M defines, by passage to the quotient, a continuous bijection \(\varphi_{x}\) of \(G/H_{x}\) onto M. Then \(\varphi_{x}\) is a homeomorphism of \(G/H_{x}\) onto M (in other words (GT, III, §2, No.~5) M is a topological homogeneous space).

Let \(x_0 \in \mathbf{M}\) . It suffices to prove (loc. cit., Prop. 15) that the mapping \(s \mapsto sx_0\) transforms every neighborhood \(\mathbf{V}\) of \(e\) in \(\mathbf{G}\) into a neighborhood of \(x_0\) in \(\mathbf{M}\) . Let \(\mathbf{W}\) be a symmetric compact neighborhood of \(e\) such that \(\mathbf{W}^2 \subset \mathbf{V}\) . By hypothesis, \(\mathbf{G}\) is the union of a sequence of compact sets, hence of a sequence \((s_n\mathbf{W})\) of translates of \(\mathbf{W}\) . Then \(\mathbf{M}\) is the union of the sequence of compact sets \((s_n\mathbf{W}x_0)\) . Since \(\mathbf{M}\) is a Baire space, there exists an index \(n\) such that \(s_n\mathbf{W}x_0\) has an interior point \(s_nwx_0\) (where \(w \in \mathbf{W}\) ). Consequently \(x_0\) is an interior point of

\[
w ^ {- 1} s _ {n} ^ {- 1} (s _ {n} \mathrm{W} x _ {0}) = w ^ {- 1} \mathrm{W} x _ {0} \subset \mathrm{V} x _ {0},
\]

so that \(Vx_{0}\) is a neighborhood of \(x_{0}\) in M.

